Die Durchschnittsgeschwindigkeit eines Intervalls ist der Mittelwert aus Anfangs- und Endgeschwindigkeit. In den ersten \tMaO\ ist die Anfangsgeschwindigkeit null, also ist die Geschwindigkeit nach $t_1 = \tMaO$
\begin{align}
 v_1 &= \vMbF = 2\cdot\vMa = \vMb
\end{align}
Die Geschwindigkeit wächst proportional zur Zeit. Nach $t_2 = 2\,t_1$ ist sie deshalb doppelt so gross:
\begin{align}
 v_2 &= \vMcF = 2\cdot\vMb = \vMc
\end{align}
Die Durchschnittsgeschwindigkeit in den nächsten \tMaO\ ist
\begin{align}
 \bar v_2 &= \vMdF = \frac{\vMb+\vMc}{2} = \vMd \approx \result{\vMdP{0}{2}}
\end{align}
Formal: $\bar v_2 = \frac{2\bar v_1+4\bar v_1}{2} = 3\,\bar v_1$, unabhängig von der Länge der beiden Intervalle, solange sie gleich lang sind.
